\section{总结与延伸}

两场讲座最终汇合成一套研究纪律：不要把能力、架构或趋势神秘化；先读数据和图，拆出隐藏任务与假设，再沿规模轴寻找交叉点。Jason 提供模型行为的微观分解，Hyung Won 提供架构演化的宏观分解；二者都要求研究者区分观察、解释与外推。

\subsection{十二条核心结论}

本节把长讲义压缩成可复用检查表。每条结论都保留适用条件，避免把 2024 年课堂判断升级成没有边界的 2026 年事实。

\begin{enumerate}
  \item 手工检查数据是在训练研究者自己的任务表示，不能被 benchmark dashboard 完全替代。
  \item Next-token prediction 通过真实 continuation 汇聚海量 latent-task supervision。
  \item 总体 cross-entropy 是不同 token 分布与任务难度的加权混合。
  \item 总体 loss 平滑不要求每个 benchmark metric 平滑。
  \item Emergent-looking curve 可能同时包含真实能力变化、阈值评分与稀疏采样。
  \item Inverse/U-shaped scaling 常提示竞争策略，应拆成子任务测量。
  \item 单点改进不能说明方法是否饱和、平移曲线还是改变斜率。
  \item 更便宜的有效计算是强大历史趋势，但不是无限保证。
  \item Inductive bias 在当前 regime 提高效率，也可能在更大规模成为瓶颈。
  \item Encoder-decoder 与 decoder-only 的差异可拆成四个局部、可消融假设。
  \item Translation 与长输入短目标解释了某些结构为什么曾经有效。
  \item 下一处需要重新审视的结构可能位于 objective、数据接口或系统交互，而不只在 architecture block。
\end{enumerate}

\subsection{一张表串起两位讲者}

为了把微观行为与宏观架构放在同一框架中，下面按“观察对象、分解单元、规模轴、风险与行动”对齐两场讲座。表格不是新结论，而是帮助读者迁移方法。

\begin{table}[H]
\centering
\begin{tabularx}{\textwidth}{L{0.18\textwidth} X X}
\toprule
维度 & Jason Wei & Hyung Won Chung \\
\midrule
观察对象 & token、prompt、task metric、loss curve & 架构、历史任务、compute trend、system cost \\
分解单元 & latent task、subskill、metric threshold & parameter sharing、layer connection、attention mask、objective \\
规模轴 & model/data/compute budget & compute availability、network depth、interaction length \\
主要风险 & 把聚合分数当机制，把突变当神秘能力 & 把历史捷径永久化，把旧 benchmark 当未来任务 \\
推荐行动 & 看数据、拆任务、画多点 scaling curve & 找主导力量、记录偏置、做减法 ablation \\
\bottomrule
\end{tabularx}
\caption{两场讲座在同一证据驱动框架中的对应关系。}
\end{table}

\subsection{实践用研究模板}

读者可以把课程方法直接用于一个新想法。先明确哪一项观察是事实，再写出至少两个竞争解释；随后选择能区分解释的预算点、子任务或架构消融。若结果只在单点成立，就不应声称改变了 scaling trajectory。

\begin{importantbox}{从课堂到实验的七步模板}
\begin{enumerate}
  \item 抽样并人工检查真实训练或评估数据；
  \item 定义连续 metric 与离散用户任务，避免只保留一个分数；
  \item 列出潜在任务、策略和结构假设；
  \item 在多个预算点运行 baseline 与 intervention；
  \item 对关键结构做单变量 ablation；
  \item 报告拟合区间、误差、失败案例和数据分布；
  \item 写清结论是观察、机制假设、个人判断还是外推。
\end{enumerate}
\end{importantbox}

\subsection{自测问题}

这些问题用于检查是否真正掌握讲座的推理，而不是只记住结论。回答时应引用公式、图中变量或反事实实验。

\begin{enumerate}
  \item 为什么 next-token objective 可以看成 massive multi-task learning？它又为什么不能保证均衡学会所有任务？
  \item 总体 loss 平滑下降时，exact-match accuracy 为什么可能突然上升？
  \item 如何用 repeat、quote repair 与 instruction following 解释 U-shaped scaling？
  \item 为什么一个研究方法在单个预算点领先，仍不足以声称更 scalable？
  \item Inductive bias 在什么条件下是帮助，在什么条件下会成为瓶颈？
  \item Encoder-decoder 到 decoder-only 的四步变换分别消除了什么假设？
  \item 为什么 long-input/short-target 数据可能特别适合 encoder-decoder？
  \item Bidirectional attention 为什么会增加多轮对话的重复计算？
  \item MLE 对开放式生成施加了什么结构假设？RLHF 又没有解决什么问题？
  \item 如果未来 compute trend 因能源约束改变，Hyung Won 的分析方法应怎样更新？
\end{enumerate}

\subsection{拓展阅读}

本节只列与课堂论证直接相关的原始材料。阅读时建议复现图的坐标和假设，而不是只摘结论。

\begin{itemize}
  \item Kaplan et al., \emph{Scaling Laws for Neural Language Models}：总体 loss 的幂律经验关系。
  \item Wei et al., \emph{Emergent Abilities of Large Language Models}：任务级曲线与 emergent behavior。
  \item Schaeffer et al., \emph{Are Emergent Abilities of Large Language Models a Mirage?}：metric choice 对曲线形状的影响。
  \item Sutton, \emph{The Bitter Lesson}：一般方法、计算与人工结构的长期关系。
  \item Vaswani et al., \emph{Attention Is All You Need}：原始 encoder-decoder Transformer。
  \item Chung et al., \emph{Scaling Instruction-Finetuned Language Models}：FLAN 系列的 instruction tuning 证据。
\end{itemize}

\end{document}
